The results of both the HAL Generalist Agent and the HF-ODR are directly cited from HAL's website \footnote{\url{https://hal.cs.princeton.edu/gaia}}, including GAIA scores and total costs. OAgent's result is from the original paper \citep{zhu2025oagents} while OWL's result is reproduced by \citet{wang2025efficient}. We also directly cite the results of AutoAgent and Magnetic-1 from \citet{chen2023autoagents}.

\subsection{Details of the Agent Frameworks}
In this section, we expand on the implementation details of other open-source complex agent frameworks. By analyzing the engineering details, we further claim that our proposed method achieves higher performance even with limited computing resources, validating the effectiveness of the proposed reflection mechanism.

\begin{itemize}
    \item \textbf{HAL Generalist Agent:} It is a single-agent framework that is built upon Smolagents' CodeAgent class with 200 max action steps and a fixed interval planning update of every 4 steps. The agent has access to the GoogleSearchTool, the default VisitWebpageTool and PythonInterpreterTool, TextInspectorTool, execute bash, edit file, file content search, and a visual inspector tool powered by GPT-4o.
    \item \textbf{HF-ODR:} It is a multi-agent framework that consists of a manager agent and a text browser agent, which is specialized in browsing webpages. The browser agent is allowed to execute up to 20 action steps with a fixed interval of planning at every 4 steps. It is also equipped with comprehensive tools to browse webpages, including GoogleSearchTool, VisitTool, page scrolling (PageUpTool and PageDownTool), FinderTool (i.e., Ctrl+F), FindNextTool (i.e., finding the next match in a Ctrl+F search), ArchiveSearchTool (i.e., search the Wayback Machine), and TextInsepctorTool. The manager agent, equipped with a text inspector tool and a visual inspector tool, has 12 maximum action steps and a plan every 4 steps.
    \item \textbf{OAgents:} It has a similar architecture to HAL Generalist Agent, where a single CodeAgent handles everything. Differently, it employs the optimal modules that are empirically validated in their experiments. The system supports up to 100 maximum action steps. Its core architecture features a multi-faceted planning module that supports subtask decomposition, staged plan calibration, and periodic plan reviews every 1 step. For factual grounding, it utilizes an ensembled toolkit comprising a multi-source search engine (Google, Wikipedia, Baidu, Bing, DuckDuckGo) alongside specialized multimodal inspectors for text, audio (Whisper-1), and visual analysis. The framework further incorporates a hierarchical memory system (summarization, vectorized retrieval, and long-term memory fusion) and a Test-Time Scaling (TTS) module that employs Best-of-N sampling and list-wise verification to enhance decision-making robustness.
    \item \textbf{Magentic-One:} It features an orchestrator agent that implements an outer loop to manage the task ledger and an inner loop to manage the progress ledger. Except for the orchestrator, the system consists of a \textit{WebSurfer} to manage web browser, a \textit{FileSurfer} to read local files of most types, a \textit{Coder} to write code, and a \textit{ComputerTerminal} to provide access to a console shell where the Coder’s programs can be executed, and where new programming libraries can be installed.
    \item \textbf{AutoAgent:} This framework utilizes a fully automated multi-agent architecture where a central LLM-powered actionable engine orchestrates specialized worker agents, including Web, Code, and File agents. It employs a Manager-Worker inspired by Magnetic-One, supporting dynamic task decomposition and sequential reasoning. The self-play agent customization module enables autonomous generation and refinement of tools, agent profiles, and workflows with natural language. The framework also integrates an agentic RAG system to ingest diverse file formats such as PDF and images.
    \item \textbf{OWL:} OWL is a multi-agent system built upon Camel with a Role-Playing architecture. Two agents, respectively, play the roles of user and assistant to solve user input tasks. The assistant agent manages various tools and specialized agents, including a \textit{Browser} system that is composed of a browser agent and a planning agent, \textit{multi-modality handling} tools including image, video, and audio, \textit{CodeExecution} tool, \textit{Excel} tool, \textit{Search} tool, and \textit{FileWrite} tool. The user agent decomposes the task and prompts the assistant agent to finish a sub-task while the assistant agent assigns specialized agents to execute and provides execution results or failures to the user agent.
\end{itemize}